\begin{solapartat}
La relació entre la longitud d'ona dels harmònics d'una corda i la longitud d'aquesta ve donada per l'expressió:
\[ n\,\frac{\lambda}{2} = L\ (n = 1, 2, 3...) \ \text{\punts{0,2}} \Rightarrow \]
% DUBTE: a l'original diu «node fonamental» (errata per «mode fonamental»); es transcriu tal com és.
El node fonamental el tindrem per $n = 1$, per tant: $\lambda = 2\,L = 64\mathrm{cm}$ \punts{0,2}. Els ventres estaran just al mig i els nodes un a cada extrem \punts{0,2}

La velocitat de propagació serà:
\[ v_p = \lambda\,\nu \ \text{\punts{0,2}} = 0,64\ \mathrm{m}\ 196\ \mathrm{Hz} = 125\mathrm{m/s} \ \text{\punts{0,2}} \]
\end{solapartat}

\begin{solapartat}
Tercer harmònic:

\figura[0.3]{sol_fig1.png}
\aladreta{\punts{0,3}}
\[ \lambda_3 = \frac{2}{3}\,L \ \text{\punts{0,1}} = 21,3\ \mathrm{cm} \Rightarrow \nu_3 = \frac{v_p}{\lambda_3} = 587\ \mathrm{Hz} \ \text{\punts{0,1}} \]
Cinqué harmònic:

\figura[0.3]{sol_fig2.png}
\aladreta{\punts{0,3}}
\[ \lambda_5 = \frac{2}{5}\,L \ \text{\punts{0,1}} = 12,8\ \mathrm{cm} \Rightarrow \nu_5 = \frac{v_p}{\lambda_5} = 977\ \mathrm{Hz} \ \text{\punts{0,1}} \]
\end{solapartat}
